% ═══════════════════════════════════════════════════════════════
%  Chapter 5 — Results and Evaluation
% ═══════════════════════════════════════════════════════════════
%
%  NOTE (read before compiling):
%  All quantitative results (match rates, object counts, status tallies,
%  sync-log totals, pipeline timings) are MEASURED values to be supplied by
%  the author. They appear as [TODO: ...] slots inside the tables and prose.
%  Screenshots are \includegraphics placeholders with "% TODO export" notes.
%  Nothing quantitative is fabricated.
%
\chapter{Results and Evaluation}
\label{chap:results}

\section*{Chapter Introduction}

Having designed and implemented the solution, this chapter evaluates its outcomes. The evaluation pursues four questions, in order: how severe was the environment drift when measured objectively by the comparator; to what extent did the synchronisation mechanism reduce that drift; how does the \gls{cicd} pipeline perform in governing changes across the environment lifecycle; and, overall, to what degree were the objectives set out in Chapter~1 achieved. The chapter closes with a discussion of the results and an honest account of the limitations of the approach.

The quantitative results presented here are drawn from recorded runs of the comparator and the synchronisation task. Throughout the chapter, measured figures are presented in the form of result tables whose values are populated from those runs.

% ─────────────────────────────────────────────────────────────
\section{Evaluation Methodology}
\label{sec:eval-method}

The evaluation rests on three independent sources of evidence, chosen so that the assessment does not depend on any single instrument:

\begin{enumerate}
  \item \textbf{Comparator measurements.} The read-only comparator was run on the two relevant environment pairs---\texttt{DEV} versus \texttt{PROD} and \texttt{UAT} versus \texttt{PROD}---both \emph{before} and \emph{after} synchronisation, providing a quantitative, repeatable measure of alignment.
  \item \textbf{Synchronisation logs.} The log table populated by the daily task records, per table, the action taken on each run, providing direct evidence of what the synchronisation did.
  \item \textbf{Pipeline execution records.} The GitHub Actions run history and the associated ServiceNow change requests provide evidence of the pipeline's governance behaviour.
\end{enumerate}

The principal metric is the \textbf{match rate}: the proportion of in-scope objects classified as \texttt{MATCH} (present on both sides and identical across all compared attributes), expressed as a percentage. Supporting metrics are the counts of objects in each status (\texttt{MATCH}, \texttt{MISSING\_IN\_SOURCE}, \texttt{MISSING\_IN\_TARGET}, \texttt{DRIFT}) and their breakdown by object type. Because the comparator persists every run, these metrics are reproducible and can be re-measured at any time to track the platform's alignment over its lifecycle.

% ─────────────────────────────────────────────────────────────
\section{Audit Results: Quantifying the Initial Drift}
\label{sec:results-audit}

The first comparator runs established a quantitative baseline for the environment drift that Chapter~2 had characterised qualitatively. Table~\ref{tab:baseline} summarises the baseline alignment of each lower environment against \texttt{PROD}.

\begin{table}[H]
  \centering
  \rowcolors{2}{lightgray}{white}
  \begin{tabularx}{\textwidth}{X c c c c c}
    \toprule
    \rowcolor{ibmblue!15}
    \textbf{Environment pair} & \textbf{Objects} & \textbf{Match} & \textbf{Missing} & \textbf{Drift} & \textbf{Match rate} \\
    \midrule
    \texttt{DEV} vs \texttt{PROD} & [TODO] & [TODO] & [TODO] & [TODO] & [TODO]\% \\
    \texttt{UAT} vs \texttt{PROD} & [TODO] & [TODO] & [TODO] & [TODO] & [TODO]\% \\
    \bottomrule
  \end{tabularx}
  \caption{Baseline Environment Alignment Measured by the Comparator (before synchronisation). \textbf{[TODO: populate from the recorded baseline comparison run.]}}
  \label{tab:baseline}
\end{table}

A breakdown by object type, shown in Table~\ref{tab:baseline-by-type}, locates the drift across the object categories compared by the tool, substantiating with object-level counts the discrepancy categories reported in Table~\ref{tab:discrepancy-categories}.

\begin{table}[H]
  \centering
  \rowcolors{2}{lightgray}{white}
  \begin{tabularx}{\textwidth}{l c c c c}
    \toprule
    \rowcolor{ibmblue!15}
    \textbf{Object type} & \textbf{Match} & \textbf{Missing in target} & \textbf{Drift} & \textbf{Match rate} \\
    \midrule
    Tables       & [TODO] & [TODO] & [TODO] & [TODO]\% \\
    Columns      & [TODO] & [TODO] & [TODO] & [TODO]\% \\
    Views        & [TODO] & [TODO] & [TODO] & [TODO]\% \\
    Procedures   & [TODO] & [TODO] & [TODO] & [TODO]\% \\
    Functions    & [TODO] & [TODO] & [TODO] & [TODO]\% \\
    Tasks        & [TODO] & [TODO] & [TODO] & [TODO]\% \\
    Streams      & [TODO] & [TODO] & [TODO] & [TODO]\% \\
    Stages       & [TODO] & [TODO] & [TODO] & [TODO]\% \\
    Sequences    & [TODO] & [TODO] & [TODO] & [TODO]\% \\
    Pipes        & [TODO] & [TODO] & [TODO] & [TODO]\% \\
    \bottomrule
  \end{tabularx}
  \caption{Baseline Drift by Object Type, \texttt{DEV}/\texttt{UAT} versus \texttt{PROD}. \textbf{[TODO: populate from the recorded baseline run; one such table per environment pair if desired.]}}
  \label{tab:baseline-by-type}
\end{table}

The comparator's results scorecard, reproduced in Figure~\ref{fig:scorecard}, presents these figures in the form seen by the engineering team during the audit.

\begin{figure}[H]
  \centering
  \includegraphics[width=0.9\linewidth]{comparator_results.png}
  \caption{Comparator results scorecard, showing the match rate and per-status object counts for an environment-pair comparison.}
  \label{fig:scorecard}
\end{figure}

% ─────────────────────────────────────────────────────────────
\section{Synchronisation Results: Post-Alignment State}
\label{sec:results-sync}

After the synchronisation task had run, the comparator was executed again on the same environment pairs to measure the change in alignment. Table~\ref{tab:before-after} contrasts the match rate before and after synchronisation.

\begin{table}[H]
  \centering
  \rowcolors{2}{lightgray}{white}
  \begin{tabularx}{\textwidth}{X c c c}
    \toprule
    \rowcolor{ibmblue!15}
    \textbf{Environment pair} & \textbf{Match rate (before)} & \textbf{Match rate (after)} & \textbf{Improvement} \\
    \midrule
    \texttt{DEV} vs \texttt{PROD} & [TODO]\% & [TODO]\% & [TODO] pp \\
    \texttt{UAT} vs \texttt{PROD} & [TODO]\% & [TODO]\% & [TODO] pp \\
    \bottomrule
  \end{tabularx}
  \caption{Environment Alignment Before and After Synchronisation (pp~=~percentage points). \textbf{[TODO: populate from the post-synchronisation comparison run.]}}
  \label{tab:before-after}
\end{table}

The synchronisation log provides a complementary, action-level view of the same run. Table~\ref{tab:synclog} tallies the actions recorded by the procedure across all in-scope schemas.

\begin{table}[H]
  \centering
  \rowcolors{2}{lightgray}{white}
  \begin{tabularx}{\textwidth}{l c X}
    \toprule
    \rowcolor{ibmblue!15}
    \textbf{Logged action} & \textbf{Count} & \textbf{Interpretation} \\
    \midrule
    \texttt{CLONED}            & [TODO] & Tables missing in the target, created by clone. \\
    \texttt{CTAS CREATED}      & [TODO] & Missing tables created via CTAS fallback (e.g.\ shared sources). \\
    \texttt{CLONE REPLACED}    & [TODO] & Structurally identical tables refreshed from \texttt{PROD}. \\
    \texttt{CTAS REPLACED}     & [TODO] & Identical tables refreshed via CTAS fallback. \\
    \texttt{STRUCTURE MISMATCH}& [TODO] & Tables with divergent structure---logged, not overwritten. \\
    \texttt{ERROR}             & [TODO] & Actions that failed and were recorded for review. \\
    \bottomrule
  \end{tabularx}
  \caption{Synchronisation Action Tally from the Daily Task Log. \textbf{[TODO: populate from the synchronisation log table.]}}
  \label{tab:synclog}
\end{table}

The proportion of tables falling under \texttt{STRUCTURE MISMATCH} is a particularly informative result: it isolates the residual divergence that the data-level synchronisation deliberately does \emph{not} resolve and that must instead be addressed by a governed DDL change through the \gls{cicd} pipeline. This division of labour---data refreshed automatically, structure changed only through the reviewed pipeline---is, by design, what keeps the synchronisation safe.

Figure~\ref{fig:before-after-chart} visualises the before-and-after match rates.

\begin{figure}[H]
  \centering
  \placeholder{Before/after match-rate chart}{To be generated from the measured values in Table~\ref{tab:before-after}.}
  \caption{Match rate before and after synchronisation, per environment pair.}
  \label{fig:before-after-chart}
\end{figure}

% ─────────────────────────────────────────────────────────────
\section{CI/CD Pipeline Results}
\label{sec:results-cicd}

The pipeline was evaluated by exercising it end to end: a representative change was developed on \texttt{dev}, promoted through \texttt{uat} to \texttt{main} via reviewed pull requests, and deployed to each environment by the workflow. The observed behaviour is summarised in Table~\ref{tab:pipeline-results}.

\begin{table}[H]
  \centering
  \rowcolors{2}{lightgray}{white}
  \begin{tabularx}{\textwidth}{X X}
    \toprule
    \rowcolor{ibmblue!15}
    \textbf{Aspect evaluated} & \textbf{Outcome} \\
    \midrule
    Validation gate on pull request & Invalid SQL is rejected at the \texttt{validate} job before merge. \textbf{[TODO: example/screenshot of a caught failure.]} \\
    Environment-scoped deployment    & Each branch deploys to its mapped environment (\texttt{dev}/\texttt{uat}/\texttt{main}~$\rightarrow$~\texttt{DEV}/\texttt{UAT}/\texttt{PROD}). \\
    Post-deployment verification     & \texttt{post\_deploy\_check.py} confirms the target state after each deployment. \\
    Mean deployment duration         & [TODO: measured duration per environment]. \\
    Change-management linkage        & Each \texttt{PROD} deployment is tied to a ServiceNow change request. \\
    Direct-to-production changes      & Eliminated: deployment to \texttt{PROD} occurs only via the pipeline identity. \\
    \bottomrule
  \end{tabularx}
  \caption{Observed Behaviour of the CI/CD Pipeline. \textbf{[TODO: complete the measured cells and attach run screenshots.]}}
  \label{tab:pipeline-results}
\end{table}

A representative GitHub Actions run is shown in Figure~\ref{fig:pipeline-run}.

\begin{figure}[H]
  \centering
  \includegraphics[width=0.9\linewidth]{github_actions_deploy_job.png}
  \caption{A GitHub Actions deployment run, showing the ordered steps executed by the pipeline for a promotion to a target environment.}
  \label{fig:pipeline-run}
\end{figure}

The qualitative result is the most consequential: the pipeline makes the governed path the \emph{only} path to production. The practice of applying changes directly to \texttt{PROD}---the proximate cause of the drift analysed in Chapter~2---is no longer possible within the new model, because production deployments are performed solely by the pipeline's authenticated identity after review and validation.

% ─────────────────────────────────────────────────────────────
\section{Evaluation Against Project Objectives}
\label{sec:results-objectives}

Table~\ref{tab:objective-eval} assesses the project against the seven objectives defined in Table~\ref{tab:objectives}, citing the evidence for each.

\begin{table}[H]
  \centering
  \rowcolors{2}{lightgray}{white}
  \begin{tabularx}{\textwidth}{c l c X}
    \toprule
    \rowcolor{ibmblue!15}
    \textbf{\#} & \textbf{Objective} & \textbf{Status} & \textbf{Evidence} \\
    \midrule
    O1 & Environment Audit       & Achieved & Comparator measures all object types across environment pairs (\S\ref{sec:results-audit}). \\
    O2 & Root Cause Analysis     & Achieved & Organisational and technical causes identified in Chapter~2. \\
    O3 & Synchronisation Strategy& Achieved & Share-based design (\S\ref{sec:sync}), deployed and measured (\S\ref{sec:results-sync}). \\
    O4 & CI/CD Pipeline          & Achieved & GitHub Actions workflow with validation, deployment, and verification (\S\ref{sec:impl-cicd}). \\
    O5 & GitHub Integration      & Achieved & Versioned repository with branch protection and pull-request review (\S\ref{subsec:impl-repo}). \\
    O6 & Governance Framework    & Achieved & Enforced review, pipeline-only deployment, ServiceNow change linkage (\S\ref{subsec:impl-workflow}). \\
    O7 & Documentation           & Ongoing  & This report and the in-repository documentation; handover in progress. \\
    \bottomrule
  \end{tabularx}
  \caption{Evaluation of the Project Against Its Objectives. \textbf{[TODO: adjust any status to reflect the final state at handover.]}}
  \label{tab:objective-eval}
\end{table}

% ─────────────────────────────────────────────────────────────
\section{Discussion}
\label{sec:results-discussion}

The results support three principal conclusions. First, the read-only comparator transformed the problem from one of anecdote into one of measurement: the drift that Chapter~2 could only describe qualitatively can now be quantified, located by object type, and tracked over time. This is itself a durable contribution, independent of the synchronisation, because it gives the team a permanent instrument for detecting regression.

Second, the share-based synchronisation proved well-matched to the constraint that \texttt{PROD} remain authoritative and untouched. By reading live production data through Secure Data Shares and refreshing the lower environments idempotently, it raises alignment without copying production storage or risking the source. Its deliberate restraint---refreshing data but only \emph{logging} structural divergence---means it does not mask the drift it is meant to reduce; structural changes are routed to the governed pipeline instead.

Third, the \gls{cicd} pipeline addresses the organisational root cause directly. The original problem was not merely technical but behavioural: in the absence of a governed path, engineers defaulted to direct production edits. By making the reviewed pipeline the only route to production, the new model removes the incentive structure that produced the drift in the first place.

% ─────────────────────────────────────────────────────────────
\section{Limitations}
\label{sec:results-limitations}

A candid evaluation must also acknowledge the limits of the work:

\begin{itemize}
  \item \textbf{Data-level synchronisation only.} The synchronisation refreshes data and clones missing tables, but it does not reconcile structurally divergent tables automatically; such cases are logged and must be resolved through a pipeline DDL change. This is a deliberate safety choice, not an oversight, but it means full structural convergence still depends on human-authored migrations.
  \item \textbf{Scope bounded by the main DAG.} Synchronisation covers the 14 databases touched by the main task DAG (Table~\ref{tab:sync-databases}). Objects outside that backbone are not synchronised, by design, but their alignment is therefore not guaranteed.
  \item \textbf{Single-user, local audit tool.} The comparator is operated locally by one engineer rather than running as a shared service; scaling it to scheduled, multi-user operation would require further work.
  \item \textbf{Daily cadence.} The synchronisation runs once per day, so the lower environments reflect production as of the last run rather than in real time.
\end{itemize}

These limitations are addressed, where relevant, in the perspectives for future work set out in the general conclusion.

% ═══════════════════════════════════════════════════════════════
\section*{Chapter Conclusion}

This chapter has evaluated the implemented solution against objective evidence. The comparator quantified the initial drift and, re-run after synchronisation, measured the resulting improvement in environment alignment; the synchronisation log corroborated those measurements at the level of individual tables; and the \gls{cicd} pipeline was shown to govern changes across the environment lifecycle while eliminating direct production deployments. Measured against the objectives of Chapter~1, the project achieved its audit, analysis, synchronisation, pipeline, integration, and governance goals, with documentation and handover continuing at the close of the internship. The general conclusion that follows reflects on the significance of these outcomes and sets out the directions in which the platform can continue to evolve.

\newpage
